\documentclass[10pt,a4paper]{amsart}
\usepackage[margin=19mm]{geometry}
\usepackage{amsmath,amssymb,mathtools}
\usepackage[scr=rsfs,cal=euler]{mathalfa}
\usepackage{xfrac}
\usepackage[unicode,hidelinks,bookmarks=false]{hyperref}
\newcommand{\ie}{i.e.\ }
\newcommand{\cf}{cf.\ }
\newcommand{\resp}{respectively\ }
\newcommand{\Subsection}{\subsection}
\providecommand{\nobreakdash}{\nobreak\hbox{-}}
\setlength{\emergencystretch}{2em}
\multlinegap0pt
\usepackage{amssymb}
\usepackage{enumerate}
\usepackage{bm}
\usepackage{tikz}
\usepackage{mathtools}
\newtheorem{lem}{Lemma}
\newtheorem{thm}{Theorem}
\newcommand{\cF}{\mathcal{F}}
\newcommand{\ee}{\varepsilon}
\title{Lifting for N-independent sets in II1 factors}
\author{Gregory Patchell, Austin Shiner}
\date{Source: arXiv:2609.10087v1 (CC BY 4.0)}
\begin{document}
\maketitle

\noindent\emph{Source and licence.} Lifting for N-independent sets in II1 factors. Version arXiv:2609.10087v1 (CC BY 4.0), distributed by arXiv under the Creative Commons Attribution 4.0 International licence, http://creativecommons.org/licenses/by/4.0/, as recorded in the arXiv licence metadata for this exact version and shown on the abstract page https://arxiv.org/abs/2609.10087v1. Full licence notice: \texttt{licenses/arxiv-2609.10087-CC-BY-4.0.txt}.\par\smallskip

\noindent\textit{Editorial scope.} Counted unit: the article's thm thm-main-technical (source0.tex lines 387--389). The article's complete original proof of it is delivered with it (source0.tex lines 391--467, one \texttt{proof} environment of 77 lines). No mathematics is written, repaired or re-derived: the statement and the proof are the article's own lines, in the article's own notation, and no retained line is substituted or re-flowed.  Retained uncounted as the source-local context the proof needs: lines 307--315.  Nothing was omitted from any retained span, so the package carries no omission record.  No retained line contains a citation, so the wrapper carries no bibliography and no bibliography entry.  No retained line contains a figure environment, an image-inclusion command or drawing code, so no image is delivered and the package declares no asset dependency.  Editorial formatting only, in addition: an AMS \texttt{amsart} preamble that restates the article's own declarations and macros used by the retained lines, namely the environments \texttt{lem}, \texttt{thm} on separate counters, together with the standard packages those lines need.  The retained environments print in this document as Lemma 1, Theorem 1 in order of appearance, whereas the article prints the same items with the numbering of its own full document.  The complete native source of the article as distributed in the arXiv e-print for this exact version is bundled unchanged as \texttt{sources/209873/source0.tex}.
\begin{lem}\label{lem-main-technical}
    Fix $N.$ There exist constants $C_1,C_2$, and $C_3$, depending only on $N,$ with the following property. Let $(M,\tau)$ be a tracial von Neumann algebra, $x = (x_1,\ldots,x_m)\in M_{\mathrm{sa},1}^m$, and $y = (y_1,\ldots,y_n)\in M_{\mathrm{sa},1}^n$, for some $m,n\ge1.$ Set $\delta = \delta(x,y)$, $\ee= \ee(x,y),$ and $\gamma = \gamma(x,y)$. Assume that $4N|\cF_N|\ee < \delta^2 - (2|\cF_N|-1)\gamma$ and set $\lambda = \frac{4N|\cF_N|\ee}{\delta^2 - (2|\cF_N|-1)\gamma}$. Then there exists $u$ a unitary in $M$ such that 
    \begin{enumerate}
        \item $\|u-1\| \le 2\lambda$;
        \item $\delta(uxu^{-1},y) \ge \delta - C_1\lambda$;
        \item $\ee(uxu^{-1},y)\le C_2\lambda^2$;
        \item $\gamma(uxu^{-1},y) \le \gamma + C_3\lambda$.
    \end{enumerate}
\end{lem} 

\begin{thm}\label{thm-main-technical}
    There exists a constant $C_0$, depending only on $N$, with the following property. Let $(M,\tau)$ be a tracial von Neumann algebra, $x = (x_1,\ldots,x_m)\in M_{\mathrm{sa},1}^m$, and $y = (y_1,\ldots,y_n)\in M_{\mathrm{sa},1}^n$, for some $m,n\ge1.$ Set $\delta = \delta(x,y)$, $\ee= \ee(x,y),$ and $\gamma = \gamma(x,y)$. If $C_0|\cF_N|\sqrt{\ee} < \delta^2 - (2|\cF_N|-1)\gamma$ and $\delta(\delta^2-(2|\cF_N|-1)\gamma)) \ge 8C_1N|\cF_N|\ee$, then  there exists a unitary $u$ in $M$ such that $\ee(uxu^{-1},y)=0$ and $\|u-1\| \le \frac{16N|\cF_N|\ee}{\delta^2 - (2|\cF_N|-1)\gamma} \le \frac{16N}{C_0}\sqrt{\ee}$.
\end{thm}

\begin{proof}
    Let $C_1,C_2,C_3\ge1$ be as in Lemma \ref{lem-main-technical}. Set $C_0 = 8N\sqrt{C_1+C_2+C_3}$. If $\ee=0$ then take $u=1$. Henceforth assume $\ee > 0$. Since $\ee \le 1$ we note the following two inequalities:

    \begin{equation}\label{eqn-def-lambda}
        \frac{4N|\cF_N|\ee}{\delta^2 - (2|\cF_N|-1)\gamma} < \frac{4N|\cF_N|\ee}{C_0|\cF_N|\sqrt{\ee}} \le \frac{4N}{C_0} \le \frac{1}{2}
    \end{equation}
    \begin{equation}\label{eqn-def-C0}
        \frac{8N|\cF_N|C_2}{\frac{C_0^2|\cF_N|}{4N} - 8NC_1 - 2(2|\cF_N|-1)C_3} \le 1.
    \end{equation}
    We also note from the hypothesis $\delta(\delta^2-(2|\cF_N|-1)\gamma)) \ge 8C_1N|\cF_N|\ee$ that
    \begin{equation}\label{eqn-pos-delta}
        \delta - 2C_1\frac{4N|\cF_N|\ee}{\delta^2 - (2|\cF_N|-1)\gamma} \ge 0.
    \end{equation}

    Let $\lambda_0 = 1$ and $u_0 = 1\in M$. We now construct inductively sequences $(u_k)$ of unitaries in $M$ and positive real numbers $\lambda_k$ with the following properties. Set $v_k = u_ku_{k-1}\cdots u_1$ and define $\delta_k = \delta(v_kxv_k^{-1},y),$ $\ee_k = \ee(v_kxv_k^{-1},y)$, and $\gamma_k = \gamma(v_kxv_k^{-1},y)$. Then the following will hold for our choices of $u_k$ and $\lambda_k$:
    \begin{enumerate}
        \item $\delta_k \ge \delta_{k-1} - C_1\lambda_k$
        \item $\ee_k \le C_2\lambda_k^2$
        \item $\gamma_k \le \gamma_{k-1} + C_3\lambda_k$
        \item $\lambda_k\le \frac{1}{2}\lambda_{k-1}$
    \end{enumerate}

    For $k=1,$ setting $\lambda_1 =  \frac{4N|\cF_N|\ee}{\delta^2 - (2|\cF_N|-1)\gamma}$, which by (\ref{eqn-def-lambda}) satisfies $\lambda_1\le \frac12$ and thus satisfies (4). From Lemma \ref{lem-main-technical} we obtain a unitary $u_1\in M$ satisfying (1)--(3) as well. Now assume $u_1,\ldots, u_l$ and $\lambda_1,\ldots,\lambda_l$ have been constructed. Define $\lambda_{l+1} = \frac{4N|\cF_N|\ee_l}{\delta_l^2 - (2|\cF_N|-1)\gamma_l}$. We first claim that $\lambda_{l+1} \le \frac{1}{2}\lambda_l$. 

    For each $1\le k \le l,$ by (\ref{eqn-pos-delta}), (1), and (4) we have that 
    \begin{align*}
        \delta_{k-1} - C_1\lambda_k &\ge \delta_0 - C_1(\lambda_1 + \ldots + \lambda_k)\\
        &\ge \delta - 2C_1\lambda_1\\
        &\ge 0.
    \end{align*}

    Now, for each $1\le k\le l,$ by (1), since $\delta_k \le 2N$, and since $\delta_{k-1}-C_1\lambda_k\ge 0$, we have
    \begin{equation*}
        \delta_k^2 \ge (\delta_{k-1}- C_1\lambda_k)^2 \ge \delta_{k-1}^2 - 2C_1\delta_{k-1}\lambda_k \ge \delta^2_{k-1} - 4NC_1\lambda_k.
    \end{equation*}

    By the previous inequality and (3), we have
    \begin{align*}
        \delta_k^2 - (2|\cF_N|-1)\gamma_k &\ge \delta^2_{k-1} - 4NC_1\lambda_k - (2|\cF_N|-1)\gamma_k \\
        &\ge \delta^2_{k-1} - 4NC_1\lambda_k - (2|\cF_N|-1)(\gamma_{k-1} + C_3\lambda_k) \\
        &= \delta^2_{k-1}  - (2|\cF_N|-1)\gamma_{k-1} - [4NC_1 + (2|\cF_N|-1)C_3]\lambda_k
    \end{align*}

    By applying the previous inequality for all $1\le k\le l$ and item (4), we have (noting that $\delta_0=\delta$ and $\gamma_0=\gamma$)
    \begin{align*}
         \delta_l^2 - (2|\cF_N|-1)\gamma_l &\ge \delta_0^2 - (2|\cF_N|-1)\gamma_0 - [4NC_1 + (2|\cF_N|-1)C_3]\sum_{k=1}^l\lambda_k\\
         &\ge \delta^2 - (2|\cF_N|-1)\gamma - [8NC_1 + 2(2|\cF_N|-1)C_3]\lambda_1 
    \end{align*}

    Now observe that by assumption, $C_0^2|\cF_N|^2\ee < (\delta^2 - (2|\cF_N|-1)\gamma)^2$ so that 
    \begin{align*}
        \delta^2 - (2|\cF_N|-1)\gamma > \frac{C_0^2|\cF_N|^2\ee}{\delta^2 - (2|\cF_N|-1)\gamma} = \frac{C_0^2|\cF_N|}{4N}\lambda_1.
    \end{align*}

    Therefore 
    \begin{align*}
        \delta_l^2 - (2|\cF_N|-1)\gamma_l &\ge \left[\frac{C_0^2|\cF_N|}{4N} - 8NC_1 - 2(2|\cF_N|-1)C_3\right]\lambda_1 \\
        &\ge 8N|\cF_N|C_2 \lambda_1.
    \end{align*}
        Note that this last inequality follows from (\ref{eqn-def-C0}).
    Now using the previous inequality, item (2), and the fact that $\lambda_l\le \lambda_1,$ we compute that
    \begin{align*}
        \lambda_{l+1} = \frac{4N|\cF_N|\ee_l}{\delta^2_l - (2|\cF_N|-1)\gamma_l} \le \frac{4N|\cF_N|C_2\lambda_l^2}{8N|\cF_N|C_2 \lambda_1} \le \frac{\lambda_l}{2}.
    \end{align*}
    This completes the claim. 

    We may now apply Lemma \ref{lem-main-technical} to $v_lxv_l^{-1}$ and $y$ to obtain a new unitary $u_{l+1}$; items (1)--(3) are straightforward to verify. We also note that $\|u_{l+1} - 1\| \le 2\lambda_{l+1}$ by Lemma \ref{lem-main-technical}.

    We now see that the sequence $(v_k)$ is Cauchy in operator norm; indeed, $\|v_{k+1} - v_k\| \le \|u_{k+1}-1\| \le 2\lambda_{k+1} \le 2^{-k}$ by item (4). Let $v$ be the limit of the sequence $(v_k)$, so that $v$ is itself a unitary in $M$. Arguing as in Lemma \ref{lem-main-technical}, we have that $\lim_k \ee(v_kxv_k^{-1},y)= \ee(vxv^{-1},y)$. But the former are bounded by $C_2\lambda_k^2 \to 0$, so $\ee(vxv^{-1},y) = 0.$ We also note that 
    \begin{align*}
        \|v-1\| &\le \sum_{k=0}^\infty \|v_{k+1}-v_k\| \\
        &\le \sum_{k=0}^\infty 2\lambda_{k+1}\\
        &\le 4\lambda_1 \\
        &= \frac{16N|\cF_N|\ee}{\delta^2 - (2|\cF_N|-1)\gamma}, 
    \end{align*}
    proving the last assertion in the statement of the theorem.
\end{proof}

\end{document}
